
\documentclass[a4paper,10pt]{report}

%\usepackage[landscape]{geometry}
\usepackage[cm-default]{fontspec}
\setromanfont{FreeSerif}
\setsansfont{FreeSans}
\setmonofont{FreeMono}
\usepackage[utf8x]{inputenc}
\usepackage{fontspec}
\usepackage{xunicode}
\usepackage{xltxtra}
\usepackage{xgreek}
\usepackage{amsmath}
\usepackage{unicode-math}
\usepackage{ulem}
\usepackage{color}
\usepackage{verbatim}
\usepackage{nopageno}
\usepackage{graphicx}
\usepackage{textpos}
\setlength{\TPHorizModule}{1cm}
\setlength{\TPVertModule}{1cm}
%\usepackage[colorgrid,texcoord]{eso-pic}
\usepackage[outline]{contour}
\usepackage{wrapfig}
\usepackage{url}
\usepackage{color}
\usepackage{tikz}
\usepackage{fancybox,fancyhdr}
\usepackage{subfigure}
\usepackage{pstricks}
\usepackage{epsfig}
\usepackage{multicol}
\usepackage{listings}
\usepackage{enumerate}
\usepackage{hyperref}
\hypersetup{
  bookmarks=true,
  bookmarksopen=true,
  pdfborder=false,
  pdfpagemode=UseNone,
  raiselinks=true,
  pdfhighlight={/P},
  colorlinks,
  citecolor=black,
  filecolor=black,
  linkcolor=black,
  urlcolor=black
}


\usepackage{fontspec}
\usepackage{xunicode}
\usepackage{xltxtra}


% Margins
%\setlength{\textwidth}{24cm}

%\setlength{\voffset}{0in}
%\setlength{\textheight}{6.8in}
% Colors
\definecolor{cbrown}{rgb}{0.49,0.24,0.07}
\definecolor{cmpez}{rgb}{0.92,0.88,0.79}
\definecolor{cmpez2}{rgb}{0.62,0.33,0.34}
\definecolor{cwhite}{rgb}{1,1,1}
\definecolor{cblack}{rgb}{0,0,0}
\definecolor{cred}{rgb}{0.9,0.15,0.15}
\definecolor{clblue}{rgb}{0.57,0.85,0.97}
\definecolor{corange}{rgb}{0.86,0.49,0.18}
\definecolor{clorange}{rgb}{0.95,0.84,0.61}
\definecolor{cyellow}{rgb}{1,0.95,0.16}
\definecolor{cgreen}{rgb}{0.69,0.8,0.31}
\definecolor{clbrown}{rgb}{0.78,0.67,0.43}
\definecolor{cmagenta}{rgb}{0.79,0.05,0.54}
\definecolor{cgray}{rgb}{0.76,0.73,0.63}
\definecolor{bluesite}{rgb}{0.05,0.43,0.67}
\definecolor{redsite}{rgb}{0.90,0.15,0.15}
\definecolor{backsite}{rgb}{0.07,0.07,0.07}
% ----------- from Costas -----------------------------

%---------------   Main Fonts    -----------------%   

% \setmainfont[Mapping=tex-text]{Calibri}
% \setmainfont[Mapping=tex-text]{Times New Roman}
%\setmainfont[Mapping=tex-text]{Droid Serif}
%\setmainfont[Mapping=tex-text]{Cambria}
%\setmainfont[Mapping=tex-text]{cm-unicode}
% \setmainfont[Mapping=tex-text]{Gentium}
% \setmainfont[Mapping=tex-text]{GFS Didot}
% \setmainfont[Mapping=tex-text]{Comic Sans MS}
 \setmainfont[Mapping=tex-text]{Ubuntu}
% \setmainfont[Mapping=tex-text]{Myriad Pro}
% \setmainfont[Mapping=tex-text]{CMU Concrete}
%\setmainfont[Mapping=tex-text]{DejaVu Sans}
%\setmainfont[Mapping=tex-text]{KerkisSans}
%\setmainfont[Mapping=tex-text]{KerkisCaligraphic}
% \setmainfont[Mapping=tex-text]{Segoe Print}
% \setmainfont[Mapping=tex-text]{Gabriola}

%------------------------------------------------%

% ------------ Mathematics Fonts ----------------%
\setmathfont{Asana-Math.ttf}
% -----------------------------------------------%
% Misc
\author{Κασωτάκη Ε. - Σμαραγδάκης Κ.}
\title{www.math24.gr}
% ---------------- Formation --------------------%


\setlength\topmargin{-0.7cm}
\addtolength\topmargin{-\headheight}
\addtolength\topmargin{-\headsep}
\setlength\textheight{26cm}
\setlength\oddsidemargin{-0.54cm}
\setlength\evensidemargin{-0.54cm}
%\setlength\marginparwidth{1.5in}
\setlength\textwidth{17cm}

\RequirePackage[avantgarde]{quotchap}
\renewcommand\chapterheadstartvskip{\vspace*{0\baselineskip}}
\RequirePackage[calcwidth]{titlesec}
\titleformat{\section}[hang]{\bfseries}
{\Large\thesection}{12pt}{\Large}[{\titlerule[0.9pt]}]
%--------------------------------------------------%
\usepackage{draftwatermark}
\SetWatermarkText{www.math24.gr}
\SetWatermarkLightness{0.9}
\SetWatermarkScale{0.6}
%--------------------------------------------------%


\usepackage{xcolor}
\usepackage{amsthm}
\usepackage{framed}
\usepackage{parskip}

\colorlet{shadecolor}{bluesite!20}

\newtheorem*{orismos}{Ορισμός}

\newenvironment{notation}
  {\begin{shaded}\begin{theorem}}
  {\end{theorem}\end{shaded}}




% Document begins
\begin{document}
\pagestyle{fancy}
\fancyhead{}
\fancyfoot{}
\renewcommand{\headrulewidth}{0pt}
\renewcommand{\footrulewidth}{0pt}


\fancyhead[LO,LE]{
 %\textblockcolor{backsite}
 %\begin{textblock}{5}(-2,-0.55)
  %\rule{0cm}{1cm}
 %\end{textblock}
 \textblockcolor{bluesite}
 \begin{textblock}{5}(-1.5,-0.55)
  \rule{0cm}{1cm}
 \end{textblock}
 %\textblockcolor{bluesite}
 %\begin{textblock}{14}(5.5,-0.55)
 % \rule{0cm}{1cm}
 %\end{textblock}
 \begin{textblock}{0}(-1,-0.25)
 \color{cwhite} \begin{Large}www.math24.gr\end{Large}
 \end{textblock}
\begin{textblock}{0}(16,-0.4)
 \color{cwhite} \includegraphics[height=1.5cm]{math24_logo.png}
 \end{textblock}
\textblockcolor{white}
\begin{textblock}{17}(-1,28)
 \color{backsite} \begin{small}Copyright \textcopyright 2012, Κασωτάκη Ε.(ikasotaki@gmail.com) - Σμαραγδάκης Κ.(kesmarag@gmail.com)\end{small}
 \end{textblock}
}
 
\begin{shaded}
\begin{center}
\huge \textbf{Φυλλάδιο Ασκήσεων}\\
%Πρόσθεση ρητών αριθμών
\end{center} 
\textbf{Μαθηματικά Α' Γυμνασίου} \hfill \textbf{Ημερομηνία Παράδοσης : \hspace{2em} }
\subsection*{Ονοματεπώνυμο :\hfill  \hspace{5em}}
\end{shaded}
\vspace{2em}
\begin{itemize}
 \item Ορισμός δύναμης ενός αριθμού
 \item Πρώτη δύναμη ενός αριθμού
 \item Δυνάμεις με βάση $1$
 \item Δυνάμεις με βάση $10$
 \item Ανάπτυγμα ενός αριθμού με χρήση των δυνάμεων του $10$
\end{itemize}

\section*{Θεωρία - Ορισμός Δύναμης\hfill \small{}}
\begin{itemize}
 \item \textbf{Δύναμη} του $α$ στη $ν$ (ή νιοστή δύναμη του $α$) ονομάζεται το γινόμενο που έχει $ν$ παράγοντες 
       ίσους με $α$, και συμβολίζεται με  $α^{ν}$.\\
       Δηλαδή $α^{ν}=α\cdot α\cdot α \cdots α$ ($α$ είναι $ν$ φορές).\\
       \textbf{π.χ } $2^{4}=2\cdot 2\cdot 2\cdot 2 $ (το $2$ είναι $4$ φορές)\\
       \textbf{π.χ } $6^{3}=6\cdot 6\cdot 6$ (το $6$ είναι $3$ φορές)
 \item Στη δύναμη $α^{ν}$ ο αριθμός $α$ ονομάζεται βάση της δύναμης και ο αριθμός $ν$ ονομάζεται εκθέτης της δύναμης.\\
\textbf{π.χ } Η δύναμη $3^{10}$ έχει βάση τον αριθμό $3$ και εκθέτη τον αριθμό $10$
 \item Η δύναμη $α^{ν}$ ονομάζεται νιοστή δύναμη του $α$ \textbf{ή} $α$ στη $ν$.\\
\textbf{π.χ } $20^{5}$ ονομάζεται πέμπτη δύναμη του είκοσι \textbf{ή} είκοσι στην πέμπτη.
\end{itemize}


\section*{Άσκηση 1  \hfill \small{20 μονάδες}}
Να συμπληρώσετε τον παρακάτω πίνακα: 
\begin{center}
 \begin{tabular}{|c|c|c|c|c|}
\hline                   
$α^{ν}$       & \textbf{Ονομασία}      &  \textbf{Βάση}       &  \textbf{Εκθέτης}     &  \textbf{Ανάπτυγμα δύναμης}   \\
\hline 
 $2^{6}$             & Έκτη δύναμη του 2 \textbf{ή} δύο στην έκτη &   2  &            & $2\cdot2\cdot2\cdot2\cdot2\cdot2$          \\
\hline 
 $6^{2}$             &                        &                      &                         &           \\
\hline 
                     &                        &     $12$             &  2                       &           \\
\hline 
               &                       &                       &                        &    $7\cdot7\cdot7$       \\
\hline 
               & Τρίτη δύναμη του εννιά                      &                       &                        &           \\
\hline
               & Τρία στην τετάρτη                      &                       &                        &           \\
\hline
$4^{4}$        &                       &                       &                        &           \\
\hline
               &                       &                       &                        & $8\cdot8\cdot8\cdot8\cdot8$          \\
\hline
               &                       &              $18$         &    $3$                    &           \\
\hline
               &                       &                       &                        &$17\cdot17\cdot17\cdot17$         \\
\hline
\end{tabular}
\end{center}

\section*{Θεωρία   \hfill \small{}}
\begin{itemize}
 \item \textbf{Πρώτη δύναμη ενός αριθμού $α$}: \\
        Η πρώτη δύναμη ενός αριθμού $α$ είναι ο ίδιος ο αριθμός $α$, δηλαδή $α^{1}=α$\\
        \textbf{π.χ} $7^{1}=7$\\
        \textbf{π.χ} $19^{1}=19$
 \item \textbf{Δυνάμεις του $1$: }\\ 
        Οι δυνάμεις του $1$ είναι όλες ίσες με $1$, δηλαδή $1^{ν}=1$\\
       \textbf{π.χ} $1^{3}=1$ \\
       \textbf{π.χ} $1^{80}=1$ 
 \item \textbf{Υπολογισμός δυνάμεων του $10$:}\\ 
       Για να υπολογισμός μια δύναμη που έχει βάση το $10$, βάζουμε το $1$ και τόσα μηδενικά  όσο 
είναι ο αριθμός του  εκθέτη της δύναμης.  \\
       \textbf{π.χ} $10^{2}=100$ (η δύναμη $10^{2}$ έχει εκθέτη $2$ και το αποτέλεσμα έχει $2$ μηδενικά)\\
       \textbf{π.χ} $10^{4}=10000$ (η δύναμη $10^{4}$ έχει εκθέτη $4$ και το αποτέλεσμα έχει $4$ μηδενικά)
       %\textbf{π.χ} $10^{2}=100$ (η δύναμη $10^{2}$ έχει εκθέτη $2$ και το αποτέλεσμα έχει $2$ μηδενικά)
\end{itemize}







\section*{Άσκηση 2  \hfill \small{20 μονάδες}}
Να υπολογίσετε το τετράγωνο, τον κύβο και την τέταρτη δύναμη των παρακάτω αριθμών. 
(Υπόδειξη: το τετράγωνο ενός αριθμού είναι η δεύτερη δύναμή του και ο κύβος ενός αριθμού είναι η τρίτη δύναμή του.)
\begin{enumerate}[1)]
 \item $1$
 \item $2$
 \item $3$
 \item $10$ 
\end{enumerate}



\section*{Άσκηση 3  \hfill \small{}20  μονάδες}
Να υπολογίσετε τις παρακάτω δυνάμεις:
\begin{enumerate}[1)]
 \item $3^{3}$
 \item $2^{7}$
 \item $6^{1}$
 \item $10^{5}$
 \item $1^{8}$
\end{enumerate}




\section*{Άσκηση 4  \hfill \small{} 20 μονάδες}
Να γράψετε με τη μορφή των δυνάμεων τα παρακάτω γινόμενα:
\begin{enumerate}[1)]
 \item $3\cdot 3\cdot 3\cdot 3\cdot 3$
 \item $4\cdot 4\cdot 4\cdot 4\cdot 4\cdot 4$
 \item $6\cdot 6\cdot 6$
 \item $2\cdot 7\cdot 7\cdot 7$
 \item $3\cdot 3\cdot 8\cdot 8\cdot 8\cdot 8$
\end{enumerate}
 
\section*{Θεωρία - Ανάπτυγμα αριθμού σε δυνάμεις του 10 \hfill \small{}}
Για να γράψουμε το ανάπτυγμα ενός αριθμού σε δυνάμεις του $10$  ακολουθούμε την παρακάτω διαδικασία:
\begin{itemize}
 \item Πρώτα παρατηρούμε την τάξη του αριθμού, δηλαδή πόσες μονάδες, δεκάδες, εκατοντάδες, χιλιάδες, 
       δεκάδες χιλιάδες, εκατοντάδες χιλιάδες, $\cdots$ έχει ο αριθμός αυτός.
 \item Έπειτα γράφουμε τις μονάδες, δεκάδες,$\cdots$ του αριθμού σε δυνάμεις του $10$.
 \item Στη συνέχεια κάθε μία από τις δυνάμεις του $10$ που βρήκαμε την πολλαπλασιάζουμε με τον αριθμό που 
       αντιστοιχεί στη δύναμη αυτή.
 \item Τέλος, προσθέτουμε όλα τα γινόμενα που βρήκαμε στο προηγούμενο βήμα.
\end{itemize}
\textbf{Παράδειγμα}\\
Ανάπτυγμα του αριθμού 345 σε δυνάμεις του 10
\begin{itemize}
 \item Ο αριθμός $345$ έχει $5$ μονάδες, $4$ δεκάδες και $3$ εκατοντάδες.\\
       Δηλαδή $345=3\text{ εκατοντάδες}+4\text{ δεκάδες}+5\text{ μονάδες}$\\
       Άρα $345=3\cdot 100+4\cdot 10+5\cdot 1$
 \item Για τις εκατοντάδες έχουμε $100=10^{2}$ και για τις δεκάδες έχουμε $10=10^{1}$.
 \item Δηλαδή $345=3\cdot 10^{2}+4\cdot 10^{1}+5$
 \item Άρα το ανάπτυγμα του αριθμού σε δυνάμεις του $10$ είναι $3\cdot10^{2}+4\cdot 10^{1}+5$.\\
       Δηλαδή $345=3\cdot10^{2}+4\cdot 10^{1}+5$.
\end{itemize}


\section*{Άσκηση 5  \hfill \small{}20  μονάδες}
Να γράψετε τους παρακάτω αριθμούς με τη χρήση των δυνάμεων του $10$:
\begin{enumerate}[1)]
 \item $783$
 \item $5247$
\end{enumerate}







%\let\thefootnote\relax\footnotetext{Επιμέλεια : Κασωτάκη Ειρήνη (ikasotaki@gmail.com),    Σμαραγδάκης Κώστας (kesmarag@gmail.com)} 


\end{document}
